\begin{table}[h!]
    \scriptsize
    \captionsetup{width=\linewidth}
    \Caption{\label{tab:friedman_com_pli} Teste de Friedman incluindo o PLIB como referencia complementar. Legenda: AT representa ACO\_R+TS e HR representa HHO+RVNS; Metodos lista os metodos comparados; $N$ e o numero de instancias; $\chi^2_F$ e a estatistica do teste; $p$ e o valor-p. O valor-p mede a evidencia contra a hipotese nula de ranks equivalentes entre os metodos; usualmente, $p<0{,}05$ indica diferenca estatisticamente significativa.}
    \IBGEtab{}{
        \begin{tabular}{lllll}
            \toprule
            Conj. & Metodos & $N$ & $\chi^2_F$ & $p$ \\
            \midrule \midrule
            TODOS & PLIB, HR, AT & 130 & 45.312 & 1.45e-10 \\
            CUBIC & PLIB, HR, AT & 30 & 40.12 & 1.94e-09 \\
            DIMACS & PLIB, HR, AT & 50 & 42.782 & 5.13e-10 \\
            HB & PLIB, HR, AT & 50 & 40.648 & 1.49e-09 \\
            \bottomrule
        \end{tabular}
    }{
    \Fonte{Autor (2026).}
}
\end{table}
